% !TEX root = ../main.tex
% GENERADO por scripts/sync_market_appendices.py; editar el Markdown de origen.
\clearpage
\section{Certificaciones y marcos de aseguramiento: alcance y evidencia}
\label{app:certificaciones}
\begingroup\footnotesize

Revisión pública inicial: \textbf{9 de septiembre de 2026}. \textbf{43 proveedores}: 20 del conjunto inicial, 22 adicionales del registro NHS y AWS HealthScribe, incorporado el \textbf{17/\allowbreak{}09/\allowbreak{}2026}. Esta ampliación no supone una nueva revisión de las otras 42 fichas. La comparación aplica un formato y unos criterios de evidencia uniformes; el número de credenciales no constituye una medida de calidad clínica. \hyperref[app:hcis]{Integraciones con HCIS} · \hyperref[app:precios]{Precios} · \hyperref[sec:candidatos]{Evaluación comparativa}.

\textbf{Criterios de clasificación de la evidencia:} «declara» significa que el proveedor anuncia la credencial; la categoría «documentación primaria examinada» identifica certificados, informes o registros cuyo contenido ha sido consultado directamente. La verificación íntegra del emisor, la vigencia y el alcance requiere el certificado o informe completo. «No localizado» indica ausencia de evidencia identificada en las fuentes públicas consultadas; no permite concluir que la credencial no exista. Las condiciones de datos son declaraciones/\allowbreak{}documentos del proveedor, salvo indicación contraria.

\subsection*{Definición y alcance de las credenciales}

\begin{longtable}{L{\dimexpr 0.23\linewidth-2\tabcolsep\relax}L{\dimexpr 0.77\linewidth-2\tabcolsep\relax}}
\toprule
\textbf{Credencial} & \textbf{Objeto y límites de acreditación} \\
\midrule
\endfirsthead
\textbf{Credencial} & \textbf{Objeto y límites de acreditación} \\
\midrule
\endhead
\midrule
\multicolumn{2}{r}{\footnotesize Continúa en la página siguiente}
\endfoot
\bottomrule
\endlastfoot
\textbf{ENS, categoría ALTA: España} & Esquema Nacional de Seguridad: requisitos organizativos y técnicos para sistemas del sector público y proveedores incluidos en su ámbito. \textbf{ALTA es la categoría más exigente}, sujeta a auditoría; importa el sistema incluido en el certificado. No acredita por sí sola exactitud clínica ni residencia UE. \href{https://www.boe.es/eli/es/rd/2022/05/03/311}{BOE, RD 311/\allowbreak{}2022}. \\
\textbf{ISO/\allowbreak{}IEC 27001} & Certifica el sistema de gestión de seguridad de la información: riesgos, accesos, incidentes, continuidad y protección de confidencialidad, integridad y disponibilidad. El alcance puede cubrir solo parte de la empresa. \href{https://www.iso.org/standard/27001}{ISO}. \\
\textbf{ISO/\allowbreak{}IEC 27701} & Sistema de gestión de privacidad: responsabilidades y controles sobre el tratamiento de datos personales. No sustituye el análisis de cumplimiento del RGPD. \href{https://www.iso.org/standard/27701}{ISO}. \\
\textbf{ISO/\allowbreak{}IEC 27017 /\allowbreak{} 27018} & \textbf{27017:} controles de seguridad específicos de servicios cloud. \textbf{27018:} protección de datos personales cuando la nube pública actúa como encargada del tratamiento. \href{https://www.iso.org/standard/43757.html}{27017}, \href{https://www.iso.org/standard/76559.html}{27018}. \\
\textbf{ISO/\allowbreak{}IEC 20000-1} & Sistema de gestión de servicios de TI: planificación, prestación y mejora del servicio. No evalúa exactitud clínica. \href{https://www.iso.org/standard/70636.html}{ISO}. \\
\textbf{ISO/\allowbreak{}IEC 42001} & Certifica la \textbf{gestión de IA}: gobernanza, riesgos, responsabilidades y mejora continua. No acredita la ausencia de contenido no sustentado ni certifica la precisión clínica del escriba. \href{https://www.iso.org/standard/42001}{ISO}. \\
\textbf{Ethik-IA: Francia} & Sello privado de supervisión humana y uso ético de IA sanitaria. Examina, entre otros aspectos, equidad, trazabilidad, datos y robustez; no es el marcado CE ni una aprobación del AI Act. \href{https://esante.gouv.fr/sites/default/files/media\_entity/documents/guide-ia\_vf.pdf}{Guía oficial francesa, §5.6}. \\
\textbf{ISO 9001 /\allowbreak{} 13485 /\allowbreak{} 14001} & \textbf{9001:} gestión general de calidad. \textbf{13485:} gestión de calidad para productos sanitarios y requisitos regulatorios. \textbf{14001:} gestión ambiental. Ninguna sustituye la evaluación clínica del producto. \href{https://www.iso.org/standard/62085.html}{9001}, \href{https://www.iso.org/standard/59752.html}{13485}, \href{https://www.iso.org/standard/14001}{14001}. \\
\textbf{SOC 1 /\allowbreak{} SOC 2 /\allowbreak{} SOC 3; tipos I y II} & Informes de auditoría: \textbf{SOC 1}, controles relevantes para información financiera; \textbf{SOC 2}, controles de seguridad y, según alcance, disponibilidad, confidencialidad, privacidad e integridad del procesamiento. \textbf{I:} diseño en una fecha; \textbf{II:} también funcionamiento durante un periodo. \textbf{SOC 3:} informe de uso general sobre controles, con menos detalle que SOC 2. \href{https://www.aicpa-cima.com/resources/landing/system-and-organization-controls-soc-suite-of-services}{AICPA}. \\
\textbf{C5 tipo II} & Auditoría de controles de seguridad cloud según el catálogo del BSI alemán; tipo II incluye eficacia durante un periodo. El auditor emite el informe, no un certificado del BSI. \href{https://www.bsi.bund.de/SharedDocs/Downloads/EN/BSI/Publications/CloudComputing/ComplianceControlsCatalogue-Cloud\_Computing-C5.pdf?\_\_blob=publicationFile\&v=3}{BSI}. \\
\textbf{ISAE 3000} & Norma para encargos de aseguramiento no financiero. Puede respaldar una auditoría de privacidad, pero su interpretación requiere examinar criterios, alcance y periodo: el nombre solo no acredita cumplimiento íntegro del RGPD. \href{https://www.iaasb.org/publications/international-standard-assurance-engagements-isae-3000-revised-assurance-engagements-other-audits-or}{IAASB}. \\
\textbf{NEN 7510 /\allowbreak{} 7512 /\allowbreak{} 7513: Países Bajos} & \textbf{7510:} seguridad de información sanitaria; \textbf{7512:} intercambio seguro; \textbf{7513:} registro de accesos y operaciones. Aplicar las normas no equivale automáticamente a poseer certificados externos de las tres. \href{https://www.nen.nl/en/zorg-welzijn/ict-in-de-zorg/informatiebeveiliging-in-de-zorg}{NEN}. \\
\textbf{Cyber Essentials /\allowbreak{} Plus: UK} & Protección frente a ciberataques comunes. La modalidad básica usa autoevaluación verificada; \textbf{Plus añade comprobación técnica independiente}. No evalúa calidad clínica. \href{https://www.ncsc.gov.uk/cyberessentials/overview}{NCSC}. \\
\textbf{HDS: Francia} & Certificación del alojamiento de datos de salud por un organismo acreditado. Su interpretación requiere comprobar el titular y las actividades cubiertas: usar un proveedor HDS no convierte automáticamente al escriba en certificado HDS. \href{https://esante.gouv.fr/produits-services/hds}{Agencia de Salud Digital}. \\
\textbf{HITRUST} & Certificaciones de controles de ciberseguridad con distintas modalidades y profundidad. Deben identificarse certificado, sistema y vigencia; mencionar el marco no demuestra estar certificado. \href{https://hitrustalliance.net/assessments-and-certifications}{HITRUST}. \\
\textbf{TX-RAMP: Texas} & Programa público de evaluación de seguridad cloud. \textbf{Nivel 2:} datos confidenciales en sistemas de impacto moderado/\allowbreak{}alto; existe también nivel 1 y estatus provisional. No es una aprobación clínica. \href{https://dir.texas.gov/information-security/texas-risk-and-authorization-management-program-tx-ramp/tx-ramp-frequently}{Texas DIR}. \\
\textbf{CE bajo MDR /\allowbreak{} UKCA médico; clases I y IIa} & Marcados de conformidad para productos sanitarios de la UE/\allowbreak{}Gran Bretaña. En software de \textbf{clase I} suele bastar la declaración del fabricante; \textbf{IIa requiere evaluación de un organismo externo}. Alcance e indicación de uso son esenciales. \href{https://eur-lex.europa.eu/eli/reg/2017/745/oj}{MDR UE}, \href{https://www.gov.uk/guidance/regulating-medical-devices-in-the-uk}{MHRA}. \\
\textbf{Registro MHRA} & Inscripción para comercializar en Gran Bretaña; no equivale por sí misma a una evaluación clínica independiente o certificación de clase IIa. \href{https://www.gov.uk/guidance/regulating-medical-devices-in-the-uk}{MHRA}. \\
\textbf{Certificado de libre venta} & Acredita la comercialización de los productos identificados y facilita su exportación. Es distinto del certificado de conformidad MDR y de un estudio clínico. \href{https://www.lakemedelsverket.se/en/medical-devices/sale}{Agencia sueca del medicamento}. \\
\textbf{LVV/\allowbreak{}Valvira A1: Finlandia} & Clasificación de sistemas de información sociosanitaria, con evaluación de seguridad recogida en Astori. \textbf{A1 no es una clase de producto sanitario MDR}; son marcos diferentes. \href{https://astori.lvv.fi/}{Registro LVV}. \\
\textbf{ISO 14971 /\allowbreak{} IEC 62366-1} & \textbf{14971:} gestión de riesgos de productos sanitarios. \textbf{62366-1:} usabilidad y prevención de errores de uso que afectan a seguridad. En las fichas son normas aplicadas, salvo certificado externo identificado. \href{https://www.iso.org/standard/72704.html}{ISO}, \href{https://webstore.iec.ch/en/publication/21863}{IEC}. \\
\textbf{DTAC /\allowbreak{} DCB0129 /\allowbreak{} DCB0160 /\allowbreak{} DSPT: NHS} & \textbf{DTAC:} criterios para evaluar tecnología digital al adquirirla. \textbf{DCB0129:} gestión de riesgos clínicos por el fabricante; \textbf{DCB0160}, por la organización que lo implanta. \textbf{DSPT:} autoevaluación de seguridad/\allowbreak{}protección de datos. «NHS compliant» no constituye un certificado externo único. \href{https://digital.nhs.uk/services/digital-technology-assessment-criteria-dtac}{DTAC}, \href{https://digital.nhs.uk/data-and-information/information-standards/governance/latest-activity/standards-and-collections/dcb0129-clinical-risk-management-its-application-in-the-manufacture-of-health-it-systems/}{DCB0129}, \href{https://www.dsptoolkit.nhs.uk/}{DSPT}. \\
\textbf{TGV /\allowbreak{} EFVP: Québec} & \textbf{TGV} (\emph{Trousse globale de vérification}): certificación del producto exigida por el sistema sanitario provincial. \textbf{EFVP:} evaluación del impacto en privacidad de su implantación, responsabilidad de la organización compradora; son requisitos distintos. \href{https://www.quebec.ca/sante/professionnels/programme-transcription-intelligence-artificielle}{Gobierno de Québec}. \\
\textbf{FCC /\allowbreak{} ISED /\allowbreak{} conformidad radio UE-UK; ensayo SGS} & En Heidi Remote, credenciales de radio del \textbf{hardware} y un ensayo antibacteriano concreto de SGS. No certifican el escriba ni convierten el micrófono en producto sanitario. \href{https://www.heidihealth.com/en-ca/docs/remote-user-manual}{Manual y referencias regulatorias}. \\
\textbf{WEEE /\allowbreak{} RoHS /\allowbreak{} EMC: hardware} & \textbf{WEEE:} gestión de residuos electrónicos; aquí, adhesión a un sistema de recogida/\allowbreak{}reciclaje. \textbf{RoHS:} restricción de sustancias peligrosas. \textbf{EMC:} compatibilidad electromagnética. Las dos últimas figuran en la declaración del fabricante del micrófono Tandem. \href{https://environment.ec.europa.eu/topics/waste-and-recycling/rohs-directive\_en}{WEEE/\allowbreak{}RoHS}, \href{https://trust.tandemhealth.ai/doc/trust?rid=6a78f0c64c2eeb0a116d6fb1\&r=e0qkrepqo674el4jxq28px\&view=true}{declaración}. \\
\textbf{AAPC /\allowbreak{} ACR-VPAT} & \textbf{AAPC:} aquí se refiere a una evaluación de codificación anunciada por Ambience. \textbf{ACR/\allowbreak{}VPAT:} informe sobre accesibilidad de software, publicado por Oracle. No son certificaciones generales de seguridad clínica. \href{https://www.ambiencehealthcare.com/}{Ambience}, \href{https://docs.oracle.com/en/corporate/accessibility/templates/t2-16521.html}{Oracle}. \\
\textbf{ISO 22301} & Gestión de continuidad de negocio: preparación, respuesta y recuperación ante interrupciones. No acredita exactitud clínica. \href{https://www.iso.org/standard/75106.html}{ISO}. \\
\textbf{Dubai AI Seal} & DCAI evalúa y clasifica empresas proveedoras de IA; el sello tiene número y categoría. No es aprobación de un producto sanitario. \href{https://dub.ai/en/ai-seal/}{DCAI}. \\
\textbf{CSA STAR niveles 1–2 /\allowbreak{} GovRAMP /\allowbreak{} IRAP} & \textbf{STAR 1:} autoevaluación de seguridad cloud. \textbf{STAR 2:} certificación o atestación con evaluación independiente; es necesario identificar modalidad y alcance. \textbf{GovRAMP:} programa de seguridad cloud con distintos niveles de verificación; es necesario verificar el estado específico. \textbf{IRAP:} evaluación frente a requisitos australianos, no certificación o acreditación del sistema. \href{https://cloudsecurityalliance.org/star/}{CSA}, \href{https://govramp.org/security-program}{GovRAMP}, \href{https://www.cyber.gov.au/business-government/protecting-devices-systems/assessment-evaluation-programs/irap}{IRAP}. \\
\textbf{IM1: NHS} & Proceso de aseguramiento de interfaces con sistemas de atención primaria, actualmente EMIS Web y SystmOne. No evalúa los modelos de IA ni acredita integración con HCIS. \href{https://digital.nhs.uk/services/digital-services-for-integrated-care/im1-pairing-integration}{NHS}. \\
\textbf{CREST: pentests} & Acredita empresas de pruebas de ciberseguridad y certifica profesionales. Un pentest realizado por ellos evalúa un alcance y una fecha; no convierte al escriba en «certificado CREST». \href{https://www.crest-approved.org/about-us/what-we-do/}{CREST}, \href{https://www.crest-approved.org/skills-certifications-careers/crest-certifications/}{certificaciones profesionales}. \\
\end{longtable}

\textbf{Evidencia contractual, técnica y clínica complementaria:} declarar cumplimiento de RGPD, HIPAA o PIPEDA no demuestra una certificación externa; HHS no reconoce los sellos privados de «certificación HIPAA». DPA/\allowbreak{}BAA son contratos, PIA/\allowbreak{}EFVP son evaluaciones de privacidad; pentests, model cards y publicaciones son evidencias complementarias, con el alcance descrito en cada ficha. \href{https://www.hhs.gov/hipaa/for-professionals/faq/2003/are-we-required-to-certify-our-organizations-compliance-with-the-standards/index.html}{HHS}. Los estudios se separan y se señala la participación del proveedor. Una insignia «Verified» del portal no sustituye al certificado. El alcance de la revisión se limita a fuentes públicas, sin acceso a expedientes privados ni confirmación directa de los proveedores.

\endgroup
